
\section{A migrating bed hump under bedload}

The classic Exner test of Hudson and Sweby~\cite{HudsonSweby2003}: a
frictionless channel 1000~m long carries a uniform subcritical flow,
10~m deep at 1~m/s ($q = 10$~m$^2$/s, $Fr = 0.1$), held by Dirichlet
boundaries at both ends, over a flat bed at $z = 0.1$~m with a hump
$z = 0.1 + \sin^2(\pi(x-300)/200)$ on $300 \le x \le 500$~m. Bedload
follows Grass's law $q_b = A_g u^3$ with $A_g = 0.01$, the bed moves by its
divergence with porosity $\lambda = 0.4$, and the inflow and outflow are
open to bedload so that the flat reaches at either end are steady. Grass is
total load, so nothing is in suspension.

The hump is low against the depth and the Froude number small, so the free
surface and the discharge barely change over it; holding both fixed,
$u = q/(w - z)$ and Exner reduces to a scalar conservation law in $z$
whose solution is carried along characteristics,
\begin{equation}
z\big(x_0 + c(z_0(x_0))\, t,\; t\big) = z_0(x_0), \qquad
c(z) = \frac{3 A_g q^3}{(1-\lambda)(w - z)^4} ,
\end{equation}
faster the higher the bed: the hump migrates downstream ($c = 8$~mm/s at
the crest, 5~mm/s on the flat) and its front steepens until the crest
overtakes the toe and a shock forms, at $t = 22\,600$~s here. The run
stops at half that time, while the solution is smooth. The frozen free
surface is the only approximation in the reference; it dips by
1~cm over the hump against a depth of 9~m and enters $c$ through
$(w-z)^4$, so the reference is good to about half a percent.

The case checks the Grass law, the bedload divergence and its bed update,
the open bedload boundaries (without them the inflow cell exports bedload
it never receives and digs a hole that travels downstream), and that the
bed volume is conserved: the flux in at the inflow equals the flux out at
the outflow over the same flat bed. The bed is compared with the
characteristic solution cell by cell at the end of the run, as an $L^1$
error relative to the hump volume, together with the crest position.

The bedload edge flux is first order (centred with Rusanov dissipation
scaled by the bed-wave speed), so the steepening front is smoothed: the
$L^1$ error is 9.9\% of the hump volume at 10~m cells and 5.5\% at the 5~m
cells used here, with the crest within a cell of the reference. This case
is what showed that the earlier bare centred flux was unstable, growing a
scour hole at the upstream toe and an overshoot at the crest.

\begin{figure}
\begin{center}
\includegraphics[width=0.9\textwidth]{bed_plot.png}
\end{center}
\caption{Bed elevation at the start, the middle and the end of the run,
every cell, against the characteristic solution.}
\end{figure}

\begin{figure}
\begin{center}
\includegraphics[width=0.9\textwidth]{bed_error_plot.png}
\end{center}
\caption{Bed error at the end of the run.}
\end{figure}

\begin{figure}
\begin{center}
\includegraphics[width=0.9\textwidth]{stage_plot.png}
\end{center}
\caption{Free surface at the end of the run.}
\end{figure}

\endinput
